% Paper-ready flowchart for the brain-metastasis discrepancy-hunter agent.
%
% Compile with:   pdflatex framework_flowchart.tex
% (or upload to Overleaf as a standalone document — no extra .sty files needed).
%
% Output is a tightly-cropped PDF with sans-serif Helvetica labels.
% Resize for a paper by changing the global `font=\small` and the node
% `text width=` values in the styles below.

\documentclass[border=8pt]{standalone}

\usepackage{tikz}
\usepackage{helvet}
\usepackage[T1]{fontenc}
\renewcommand{\familydefault}{\sfdefault}

\usetikzlibrary{positioning, shapes.geometric, arrows.meta, fit, backgrounds, calc}

% ---- palette (muted, paper-friendly) -----------------------------------
\definecolor{Cinput}   {HTML}{2E4B6B}
\definecolor{Cdetect}  {HTML}{8E4B2E}
\definecolor{Cagent}   {HTML}{3F6B43}
\definecolor{Coutput}  {HTML}{6B3F6B}
\definecolor{Cside}    {HTML}{555555}

\begin{document}
\begin{tikzpicture}[
    font=\small,
    node distance = 7mm and 12mm,
    >=Stealth,
    every node/.style = {align=center},
    %
    io/.style = {
      draw=Cinput, line width=0.7pt, fill=Cinput!8, rounded corners=2pt,
      minimum height=11mm, text width=42mm, inner sep=4pt,
    },
    det/.style = {
      draw=Cdetect, line width=0.7pt, fill=Cdetect!8, rounded corners=2pt,
      minimum height=11mm, text width=52mm, inner sep=4pt,
    },
    ag/.style = {
      draw=Cagent, line width=0.7pt, fill=Cagent!8, rounded corners=2pt,
      minimum height=11mm, text width=58mm, inner sep=4pt,
    },
    agdim/.style = {
      draw=black!30, line width=0.5pt, fill=black!4, rounded corners=2pt,
      minimum height=10mm, text width=58mm, inner sep=4pt,
      text=black!55,
    },
    out/.style = {
      draw=Coutput, line width=0.7pt, fill=Coutput!8, rounded corners=2pt,
      minimum height=11mm, text width=52mm, inner sep=4pt,
    },
    sidebox/.style = {
      draw=Cside, line width=0.6pt, fill=Cside!5, rounded corners=2pt,
      text width=46mm, inner sep=5pt, font=\footnotesize, align=left,
    },
    arr/.style    = {->, line width=0.6pt, draw=black!75},
    arrdash/.style= {->, dashed, line width=0.5pt, draw=black!50},
    %
    lane/.style = {
      font=\bfseries\footnotesize\sffamily, rotate=90, anchor=south,
      inner sep=2pt,
    },
]

% ============================================================
% INPUT
% ============================================================
\node[io] (mri) {%
  \textbf{4-channel MRI study}\\[1pt]
  \footnotesize T1pre\,$\cdot$\,T1post\,$\cdot$\,FLAIR\,$\cdot$\,Subtraction%
};
\node[io, right=14mm of mri] (seed) {%
  \textbf{Seed lesion mask}\\[1pt]
  \footnotesize first finding (binary)%
};

% ============================================================
% DETECTOR
% ============================================================
\node[det, below=14mm of mri.south, anchor=north,
      xshift=12mm] (nnu) {%
  \textbf{nnU-Net Task115\_Metastases\_All}\\
  \footnotesize Rudie et al.\ 2024 \,$\cdot$\, frozen \,$\cdot$\, 4-channel%
};
\node[det, below=of nnu] (cc) {%
  \textbf{Connected components $\to$ candidate pool}\\
  \footnotesize size-ranked probmap\;
  $p_i \!=\! 0.3 + 0.7\,s_i / s_{\max}$,\,
  $\geq\!3$ vox%
};

% ============================================================
% AGENT  (rule-based 7-step orchestrator)
% ============================================================
\node[ag, below=14mm of cc.south, anchor=north] (s1) {%
  \textbf{Step 1.\ characterize\_seed}\\
  \footnotesize $\to$ \textsc{SeedPhenotype}: vol, $\Delta$enh, eccentricity,
  modality\,$z$%
};
\node[ag, below=of s1] (s2) {%
  \textbf{Step 2.\ per-candidate verify} \;\textit{(loop, budget 32)}\\
  \footnotesize \textsc{EvidenceCard} (4 mod $\times$ 3 axes + seed strip)\\
  \footnotesize $\to$ Qwen2.5-VL $\to$ \textsc{VerifierVerdict}\\
  \footnotesize $\{$decision,\,vlm\_conf,\,seed\_sim,\,mimic\_risk$\}$%
};
\node[ag, below=of s2] (s3) {%
  \textbf{Step 3.\ residual\_update}\\
  \footnotesize subtract confirmed lesions from probmap%
};
\node[agdim, below=of s3] (s45) {%
  \textit{Step 4 (lower-threshold pass) and
   Step 5 (residual top-K)}\\[1pt]
  \scriptsize \textit{ablated to no-op in production --- see negative results
  \S\,6}%
};
\node[ag, below=of s45] (s6) {%
  \textbf{Step 6.\ phenotype\_compare}\\
  \footnotesize per-candidate seed-similarity (intensity\,$z$)%
};
\node[ag, below=of s6] (s7) {%
  \textbf{Step 7.\ rank\_candidates}\\
  \footnotesize \texttt{verdict\_heavy} weights\,$\cdot$\,
   prob\_source = \texttt{vlm\_else\_detector}%
};

% ============================================================
% OUTPUT
% ============================================================
\node[out, below=14mm of s7] (output) {%
  \textbf{Top-$K$ ranked additional lesions}\\
  \footnotesize + structured audit log of every step \& verdict%
};

% ============================================================
% VLM SIDECAR (training stack details)
% ============================================================
\node[sidebox, right=14mm of s2] (vlm) {%
  \textbf{Local VLM stack}\\[2pt]
  $\bullet$ Qwen2.5-VL-7B-Instruct \;\scriptsize(frozen)\\
  $\bullet$ {\small +} LoRA SFT v2 \;\scriptsize(2{,}901 ex, 2 ep)\\
  $\bullet$ {\small +} LoRA DPO v4 \;\scriptsize(103/103 pairs, $\beta\!=\!0.05$, 1 ep)\\[2pt]
  \scriptsize \textit{$r\!=\!16$, $\alpha\!=\!32$ on q/k/v/o\_proj of language tower;
   vision encoder frozen.}%
};

% ============================================================
% ARROWS
% ============================================================
% input -> detector
\draw[arr] (mri) -- (nnu);

% seed feeds the agent directly (skips detector); residual\_update also
% references the seed mask but that is a downstream consequence of
% confirming the seed at Step 1.
\draw[arr] (seed.south) |- (s1.east);

% detector chain
\draw[arr] (nnu) -- (cc);

% detector -> agent
\draw[arr] (cc) -- (s1);

% agent chain
\draw[arr] (s1) -- (s2);
\draw[arr] (s2) -- (s3);
\draw[arr] (s3) -- (s45);
\draw[arr] (s45) -- (s6);
\draw[arr] (s6) -- (s7);

% agent -> output
\draw[arr] (s7) -- (output);

% sidecar dashed link to step 2
\draw[arrdash] (vlm.west) -- (s2.east);

% ============================================================
% LANE LABELS (vertical coloured bars on the left)
% ============================================================
\begin{pgfonlayer}{background}
  % INPUT lane
  \node[fit=(mri)(seed), draw=none, fill=Cinput!4, rounded corners=2pt,
        inner xsep=2mm, inner ysep=3mm] (laneIn) {};
  % DETECTOR lane
  \node[fit=(nnu)(cc), draw=none, fill=Cdetect!4, rounded corners=2pt,
        inner xsep=2mm, inner ysep=3mm] (laneDet) {};
  % AGENT lane
  \node[fit=(s1)(s2)(s3)(s45)(s6)(s7)(vlm), draw=none, fill=Cagent!4,
        rounded corners=2pt, inner xsep=2mm, inner ysep=3mm] (laneAg) {};
  % OUTPUT lane
  \node[fit=(output), draw=none, fill=Coutput!4, rounded corners=2pt,
        inner xsep=2mm, inner ysep=3mm] (laneOut) {};
\end{pgfonlayer}

% lane labels (rotated, on the far left)
\node[lane, text=Cinput,  anchor=center] at ($(laneIn.west)+(-3mm,0)$)  {INPUT};
\node[lane, text=Cdetect, anchor=center] at ($(laneDet.west)+(-3mm,0)$) {DETECTOR};
\node[lane, text=Cagent,  anchor=center] at ($(laneAg.west)+(-3mm,0)$)  {AGENT \,(rule-based 7-step)};
\node[lane, text=Coutput, anchor=center] at ($(laneOut.west)+(-3mm,0)$) {OUTPUT};

\end{tikzpicture}
\end{document}
